\chapter{Dinámica relacional de marcos: autoinercia, ambivalencia y principio
de Mach}
\label{chap:autoinercia-mach}

La inercia clásica caracteriza la persistencia de un estado de movimiento.
La pregunta de Mach \cite{MachMechanics} desplaza el énfasis desde el cuerpo aislado hacia la
relación entre ese cuerpo y la totalidad material. En el lenguaje
genealógico, la misma cuestión se formula sin separar al observador de la
frontera que lee: un marco local es la restricción de un estado compatible
global, y la respuesta inercial es machiana cuando depende de la totalidad a
través de una traza de frontera bien definida.

La palabra \emph{autoinercial} no designa un sistema desconectado del resto.
Designa una sección que transporta coherentemente el marco inducido por su
propia inserción en la totalidad. Su persistencia es interna al estado global
y se comprueba mediante una conexión.

\section{Totalidad, traza y respuesta local}

Sea \(\Omega_{\rm tot}\) el espacio de secciones globales y
\(b:\Omega_{\rm tot}\to\mathcal B\) su traza sobre una frontera. Para una
región o vértice \(v\), sea
\[
 \mathcal I_v:\Omega_{\rm tot}\longrightarrow\mathcal Y_v
\]
la respuesta inercial local.

\begin{definition}[Factorización machiana]
La respuesta \(\mathcal I_v\) es machiana respecto de \(b\) cuando existe
un mapa inducido
\(\overline{\mathcal I}_v:\operatorname{im}(b)\to\mathcal Y_v\) tal que
\begin{equation}
 \mathcal I_v=\overline{\mathcal I}_v\circ b.
 \label{eq:factorizacion-machiana}
\end{equation}
\end{definition}

\begin{theorem}[Criterio de factorización por totalidad]
\label{thm:criterio-mach}
Existe un único mapa
\(\overline{\mathcal I}_v:\operatorname{im}(b)\to\mathcal Y_v\) que satisface
\eqref{eq:factorizacion-machiana} si y sólo si \(\mathcal I_v\) es constante
sobre cada fibra de \(b\):
\[
 b(x)=b(y)\quad\Longrightarrow\quad
 \mathcal I_v(x)=\mathcal I_v(y).
\]
Si \(b\) es sobreyectiva, \(\operatorname{im}(b)=\mathcal B\) y el mapa
factorizado queda definido de modo único sobre toda \(\mathcal B\).
\end{theorem}

\begin{proof}
La necesidad es inmediata. Para la suficiencia, dado
\(\beta\in\operatorname{im}(b)\), definimos
\(\overline{\mathcal I}_v(\beta)=\mathcal I_v(x)\) para cualquier
\(x\in b^{-1}(\beta)\). La constancia en las fibras garantiza que la
definición no depende del representante; la propia igualdad de factorización
fuerza su valor en cada \(\beta\in\operatorname{im}(b)\), y por ello garantiza
la unicidad. Si \(b\) es sobreyectiva, este dominio coincide con
\(\mathcal B\).
\end{proof}

El criterio da una formulación positiva del principio de Mach: la frontera
global contiene exactamente la información de la totalidad que gobierna la
respuesta local. El estado interior puede conservar más genealogía; la
inercia depende de la clase que la traza publica.

\section{Secciones autoinerciales como geodésicas del estado conjunto de sistema, observador y memoria: \(D_\gamma s_\gamma=0\)}

Sea \(\mathcal E\to\Gamma\) un fibrado vectorial de estados locales a lo largo
de una ruta \(\gamma\), y sea \(D\) una conexión lineal inducida por el
transporte TPK.

\begin{definition}[Sección autoinercial]
Una sección persistente \(s_\gamma\) es autoinercial en un tramo si
\begin{equation}
 D_\gamma s_\gamma=0.
 \label{eq:autoinercia-paralela}
\end{equation}
Equivalente y discretamente, el estado en cada vértice es el transporte
paralelo del estado anterior por el morfismo de la ruta.
\end{definition}

La ecuación \eqref{eq:autoinercia-paralela} conserva orientación, memoria y
firma en el sentido determinado por la conexión. Una aceleración coordenada
puede coexistir con transporte paralelo cuando el marco de referencia también
cambia. La distinción entre movimiento inercial y fuerza queda así ligada a
la conexión, no a la apariencia de la trayectoria en una carta aislada.

\begin{proposition}[Persistencia y holonomía]
\label{prop:autoinercia-holonomia}
Sea \(\gamma\) un lazo y \(\Hol_D(\gamma)\) su holonomía. Una sección
autoinercial vuelve exactamente a su estado inicial si y sólo si pertenece al
subespacio fijo de \(\Hol_D(\gamma)\). Si la proyección visible vuelve y el
estado inicial queda fuera de
\(\operatorname{Fix}(\Hol_D(\gamma))\), retorna la coordenada y cambia el
estado.
\end{proposition}

Este resultado enlaza autoinercia y monodromía. La periodicidad observable no
implica cierre del lift; el déficit o exceso de holonomía registra la
respuesta acumulada del marco.

\section{Principio de Ambivalencia: involución bidireccional de los lectores π/φ² y φ/π²}

La celda HMT admite dos lecturas complementarias: areal y volumétrica. Sea
\(q_A\) el lector de frontera o área y \(q_V\) el lector de interior o
volumen. El transporte \(T_{AV}\) relaciona ambos sin identificarlos.

\begin{principle}[Ambivalencia areal--volumétrica]
Una realización física completa conserva el par
\begin{equation}
 \bigl(q_A(x),q_V(x),T_{AV}(x)\bigr).
 \label{eq:ambivalencia-av}
\end{equation}
La lectura areal expresa cómo el estado incide en una frontera; la lectura
volumétrica expresa cómo ocupa y transporta el interior. El tercer dato
conserva la transformación entre ambas.
\end{principle}

En la Mecánica Dimensional, la relación areal--volumétrica organiza la
extensión HMT del principio de equivalencia. Una misma sección puede leerse
como respuesta de frontera y como densidad interior. Cuando ambas lecturas
son transportadas por la misma conexión local, las descripciones inercial y
gravitatoria pertenecen a una sola clase de estado.

\begin{criterion}[Equivalencia local tipada]
La identificación de una respuesta inercial con una respuesta gravitatoria
queda demostrada en una región \(U\) cuando:
\begin{enumerate}
 \item \(q_A\) y \(q_V\) se definen sobre la misma sección;
 \item existe un entrelazador \(T_{AV}\) que conserva sus unidades y
       orientaciones;
 \item la conexión que transporta ambos lectores coincide en \(U\);
 \item la diferencia observacional se anula en el marco local declarado.
\end{enumerate}
\end{criterion}

La formulación preserva la equivalencia y muestra qué datos la realizan. Su
aplicación al electrón pertenece a la primera clausura física de la ley de
masas; el presente anexo utiliza sólo el esquema geométrico.

\section{Coordinación causal entre transporte inercial, doble lectura areal–volumétrica y factorización machiana}

Los tres principios se ordenan causalmente:
\begin{equation}
 \boxed{
 \begin{aligned}
 &\text{persistencia de la sección}
   &&\longrightarrow&& D_\gamma s=0,\\
 &\text{doble lectura del mismo estado}
   &&\longrightarrow&& (q_A,q_V,T_{AV}),\\
 &\text{dependencia de la respuesta local de la totalidad}
   &&\longrightarrow&& \mathcal I_v=\overline{\mathcal I}_v\circ b.
 \end{aligned}}
 \label{eq:triptico-inercia-mach}
\end{equation}
El principio de inercia fija el transporte; la Ambivalencia fija las dos
lecturas físicas; Mach fija el cociente global que gobierna el marco local.
Juntos producen una teoría relacional del movimiento en la que el observador,
el aparato y el cuerpo pertenecen a una sección común.

\section{Medición de movimiento y acoplamiento de marcos}

Medir velocidad o aceleración exige comparar el sistema con un reloj y una
regla. Esos objetos forman parte del aparato. Sea \(F_S\) el marco del sistema
y \(F_M\) el marco del observador. El acoplamiento de medida construye una
transformación relativa
\[
 G_{MS}=F_M^{-1}F_S.
\]
La velocidad, la aceleración y la fase publicadas son invariantes o
componentes de \(G_{MS}\), no propiedades sin referencia. La repetición del
experimento conserva el contexto al transportar conjuntamente aparato y
frontera.

Esta formulación recupera la intuición de Weyl \cite{Weyl1927}: medir una amplitud o una
longitud implica acoplar una regla de calibre. HMT añade la genealogía de esa
regla. El resultado puede remontarse al estado del aparato, al lector y al
transporte que fijaron su escala.

\section{Onda bidireccional y orientación temporal}

Una ruta persistente admite los canales avanzado y retardado. La inversión
de orientación actúa sobre la palabra y el ledger; en la realización lineal,
intercambia los propagadores correspondientes cuando el Hamiltoniano posee la
simetría declarada. El movimiento bidireccional se expresa mediante un par
\[
 (U_{\rm ret},U_{\rm adv})
\]
relacionado por conjugación y reversión temporal.

La banda partícula--antipartícula realiza esta dualidad como dos orientaciones
de una misma sección. La relación dinámica puede expresarse de forma exacta.

\begin{proposition}[Conjugación de orientación y reversión del propagador]
Sea \(C\) una involución ortogonal sobre la realificación cuántica tal que
\[
 C J_E C^{-1}=-J_E,
 \qquad
 C H C^{-1}=H.
\]
Entonces, para
\(U(t)=\exp(-J_EHt/\hbar_{\rm int})\),
\[
 C U(t)C^{-1}=U(-t).
\]
\end{proposition}

\begin{proof}
La conjugación conmuta con la serie exponencial y transforma el generador
\(-J_EH\) en \(+J_EH\). El resultado es la exponencial evaluada en
\(-t\).
\end{proof}

La involución de palabra \(C_M=-\operatorname{rev}\) proporciona el dato
genealógico de orientación. Cuando su realización satisface las dos
identidades de la proposición, las ramas partícula y antipartícula se
representan mediante propagadores de sentidos temporales conjugados. La
geometría de banda fija la orientación y el Hamiltoniano fija la dinámica.

\section{Expansión global y ecuación de autoconsistencia}

Una cosmología genealógica puede formularse mediante un operador
\(\mathcal T\) que lleva una totalidad de frontera a la siguiente:
\[
 B_{n+1}=\mathcal T(B_n).
\]
La respuesta local inducida es
\[
 \mathcal I_{v,n}=\overline{\mathcal I}_v(B_n).
\]
Un régimen autoinercial global corresponde a una órbita o punto fijo estable
del operador de transporte, mientras la expansión se registra mediante un
cociclo de escala \(a_{n+1}/a_n\). Esta formulación separa tres preguntas:
existencia de la dinámica, estabilidad de la solución e identificación del
cociclo con el factor cosmológico observado.

\begin{resultbox}
La inercia local puede tiparse como transporte paralelo de una sección; su
dependencia machiana se expresa como factorización por la traza global de
frontera. El principio de Ambivalencia enlaza las lecturas areal y volumétrica
del mismo estado. El observador entra mediante el acoplamiento de marcos que
hace posible la medida.
\end{resultbox}

\begin{statusbox}
La definición de autoinercia y la factorización machiana son resultados
recuperados. El criterio de dimensión y el instrumento de medida proceden de
los capítulos anteriores. La síntesis \eqref{eq:triptico-inercia-mach} es una
formalización nueva de arquitectura autoral preexistente. Una ecuación
cosmológica cuantitativa requiere fijar \(\mathcal T\), su métrica y sus datos
de frontera.
\end{statusbox}

La atribución machiana queda caracterizada por la constancia de la respuesta
local sobre las fibras de la traza global. La autoinercia queda caracterizada
por el transporte paralelo y por el subespacio fijo de la holonomía. En una
realización cosmológica, el operador \(\mathcal T\), su métrica y el cociclo
de escala convierten esas condiciones estructurales en una evolución
cuantitativa de la totalidad.
